\documentclass[a4paper]{article}
\usepackage{amsfonts}
\usepackage{amsmath}
\usepackage{amssymb}
\usepackage{graphicx}

\usepackage{cancel}

\begin{document}
  
\noindent \large Auteur: Siebe Mees \\
\noindent \large Studierichting: Informatica\\
\noindent \large Oefeningenreeks 6E nr. 2.3\\

\medskip

\normalsize

$T_n\left(f,0\right)(x)$ voor $f(x)=\cos x$ en bewijs $\lim\limits_{n\rightarrow +\infty}R_n(x)=0$\\

\begin{enumerate}

\item $T_n\left(\cos x,0\right)(x)$\\

$f(x)=\cos x$\\

$f'(x)=-\sin x$\\

$f''(x)=-\cos x$\\

$f'''(x)=\sin x$\\

$f^{iv}(x)=\cos x$\\

$f(0)=1$\\

$f'(0)=0$\\

$f''(0)=-1$\\

$f'''(0)=0$\\

$f^{iv}(0)=1$\\

$\Rightarrow f^{(2k+1)}(0)=0 $ en $ f^{(2k)}(0)=(-1)^k$\\

$\Rightarrow T_n\left(\cos x,0\right)(x)=\displaystyle\sum_{k=0}^{\infty}\dfrac{f^{(k)}(0)}{k!}x^k$\\

Alleen even termen blijven over:\\

$=\displaystyle\sum_{k=0}^{\infty}\dfrac{f^{(2k)}(0)}{(2k)!}x^{2k}=\displaystyle\sum_{k=0}^{\infty}\dfrac{(-1)^k}{(2k)!}x^{2k}$\\

\item Bewijs $\lim\limits_{n\rightarrow +\infty}R_n(x)=0$\\

$R_n(x)=\dfrac{f^{(2n+2)}(\xi)}{(2n+2)!}x^{2n+2}\quad \xi\in[0,x]$\\

$\left|f^{(m)}(x)\right|\leq 1$\\

$\left|R_n(x)\right|\leq\dfrac{1}{(2n+2)!}|x|^{2n+2}$\\

Gebruik D'Alembert: Beschouw $\displaystyle\sum_{n=0}^{\infty}\dfrac{x^{2n}}{(2n)!}$\\

$\lim\limits_{n\rightarrow +\infty}\dfrac{a_{n+1}}{a_n}=\lim\limits_{n\rightarrow +\infty}\dfrac{x^{2n+2}}{(2n+2)!}\cdot\dfrac{(2n)!}{x^{2n}}$\\

$=\lim\limits_{n\rightarrow +\infty}\dfrac{x^2}{(2n+2)(2n+1)} = \dfrac{x^2}{\infty}=0 \Rightarrow$ De reeks is convergent.\\

Bijgevolg moet $\lim\limits_{n\rightarrow +\infty}R_n(x) = \lim\limits_{n\rightarrow +\infty}\dfrac{x^{2n+2}}{(2n+2)!}=0$  \\

\end{enumerate}

\begin{thebibliography}{999}
%\bibitem{a} Referentie
\end{thebibliography}
\end{document} 
